\chapter{Proposta de Solução}\label{chap:proposta_solucao}


O artefato proposto é um sistema multiagente composto por dois agentes com funções distintas. O \textbf{agente de conformidade de gênero} verifica se um manuscrito realiza as regras de gênero de uma trilha de conferência. O \textbf{agente de fronteira de rejeição} sinaliza os riscos de recusa a que o manuscrito está exposto, a partir de avaliações reais de artigos rejeitados. A integração das duas análises em um retorno único ao autor é a meta da versão final do artefato.


% \section{Unidade de análise}
% \label{sec:unidade}

% A Teoria de Gêneros admite descrição em diferentes níveis de resolução, e o nível adotado é uma decisão de projeto que precisa ser declarada \cite{GenreRepertoire}. Uma conferência pode ser tratada como sistema de gêneros, cujas trilhas constituem subgêneros distintos, ou como gênero único, cujas trilhas seriam apenas recortes temáticos de uma mesma ação comunicativa. As duas leituras produzem artefatos diferentes.

% Este trabalho adota a \textbf{trilha} como unidade de análise. A escolha decorre da própria fundamentação. \citeonline{GenresOrgComm} sustentam que variantes especializadas de um gênero existem sempre que a uma delas correspondam situação recorrente, propósito comum e características formais compartilhadas, e a trilha satisfaz as três condições. Uma trilha define quem submete, o que conta como contribuição pertinente e sob que forma ela é reconhecida, e faz isso de modo distinto das demais trilhas do mesmo evento.


\section{Requisitos derivados da fundamentação}
\label{sec:requisitos}

Em Design Science Research, os requisitos de um artefato derivam da base de conhecimento que justifica a intervenção, e não de levantamento junto a usuários \cite{Hevner2004}. A Tabela~\ref{tab:requisitos} registra a correspondência entre cada achado teórico e o requisito que dele decorre.

\begin{table}[H]
\centering
\caption{Derivação dos requisitos de projeto a partir da fundamentação teórica}
\label{tab:requisitos}
\small
\begin{tabular}{|p{6.4cm}|p{7.4cm}|}
\hline
\textbf{Origem teórica} & \textbf{Requisito de projeto} \\
\hline
Regras de gênero são tácitas e distribuídas nos exemplares \cite{GenresOrgComm} & A especificação deve ser derivada do corpus de artigos aceitos, e não estipulada a priori pelo pesquisador \\
\hline
O reconhecimento exige regras distintivas \textit{suficientes}, e não a totalidade delas \cite{GenresOrgComm} & A verificação deve apurar quais regras o manuscrito invoca e quais deixa de invocar, sem exigir realização integral \\
\hline
Gêneros variam por comunidade e por nível de abstração \cite{GenreRepertoire} & A especificação deve ser parametrizada por trilha, e não única para a conferência \\
\hline
O \textbf{move} é delimitado por função retórica, e não por extensão \cite{GenreAnalysis} & A unidade de segmentação e de comparação do sistema deve ser o \textbf{move} \\
\hline
Elementos são obrigatórios, opcionais ou prováveis em determinadas áreas \cite{ResearchGenres} & O retorno deve distinguir os três casos, sem tratar ausência de elemento opcional como falha \\
\hline
A dimensão do propósito permanece pressuposta na chamada \cite{GenreSystems} & O artefato deve tornar explícito o que a chamada não enuncia, complementando-a e não a repetindo \\
\hline
Usuários incorporam orientação de modelos de linguagem sem análise crítica \cite{WhyandWhen} & Toda afirmação sobre o texto deve ser rastreável, com evidência textual e exemplares verificáveis \\
\hline
Modelos classificam erroneamente falhas e sofrem influência de recomendações fortes \cite{WhereLLMWrong} & A decisão de conformidade deve ser determinística e baseada nas regras de gênero \\
\hline
\end{tabular}
\end{table}

As duas últimas linhas restringem o projeto mais do que o orientam. O achado de \citeonline{WhyandWhen}, de que usuários incorporam orientação de modelos sem submetê-la a exame, impede que o artefato produza qualquer afirmação que o autor não possa verificar por conta própria. O de \citeonline{WhereLLMWrong}, de que modelos confundem falhas de mérito com falhas de apresentação e se deixam arrastar por recomendações fortes, impede que decisões de conformidade sejam delegadas a um modelo. Esses dois requisitos valem para os dois agentes e explicam a divisão de trabalho entre componentes determinísticos e inferenciais que estrutura cada um deles.


\section{Agente de conformidade de gênero}
\label{sec:agente-conformidade}

O agente de conformidade de gênero modela o gênero comunicativo de uma trilha a partir do corpus de artigos nela aceitos e verifica a conformidade de um manuscrito em elaboração a essa especificação, produzindo retorno rastreável e formativo ao autor.

O produto entregue ao usuário são duas respostas. A primeira é a verificação de conformidade de gênero: dado um manuscrito e definida uma trilha-alvo, o agente identifica quais regras de gênero o manuscrito contempla, quais deixa de contemplar e em que ponto do texto cada função retórica esperada foi ou não realizada. A segunda é a orientação de escrita, que apresenta realizações e ausências como retorno acionável, ancorado na evidência textual do manuscrito e em exemplares reais extraídos do corpus da trilha.

A análise feita por este agente não julga o mérito científico da contribuição, não estima probabilidade de aceitação, não emite recomendação editorial, não avalia correção metodológica e não reescreve o texto do autor. Cada uma dessas exclusões corresponde a uma tarefa de natureza distinta da verificação de pertencimento a um gênero.

O agente opera onde coexistam três condições: uma comunidade científica com gêneros estáveis, um corpus de instâncias aceitas suficiente para caracterizá-los e um insumo textual estruturável.

\subsection{Corpus}
\label{sec:corpus-conformidade}

O corpus reúne artigos completos aceitos e publicados em conferencias brasileiras, para a construção deste artefato foram coletados os artigos publicados nos \textit{Anais do XXII Simpósio Brasileiro de Sistemas de Informação (SBSI 2026)} e nos \textit{Anais do XXII Encontro Nacional de Inteligência Artificial e Computacional (ENIAC 2025)}, obtidos na SBC OpenLib \cite{solsbc}. A Tabela~\ref{tab:corpus} apresenta sua organização em três conjuntos.

\begin{table}[H]
\centering
\caption{Composição do corpus por trilha}
\label{tab:corpus}
\small
\begin{tabular}{|p{4.4cm}|p{3.4cm}|p{1.6cm}|p{4.0cm}|}
\hline
\textbf{Evento} & \textbf{Trilha} & \textbf{Artigos} & \textbf{Função no estudo} \\
\hline
SBSI 2026 & Trilha de Pesquisa em SI & 67 & Corpus-alvo \\
\hline
ENIAC 2025 & \textit{Main track} & 108 & Contraste entre comunidades \\
\hline
ENIAC 2025 & \textit{Undergraduate track} & 69 & Contraste entre trilhas do mesmo evento \\
\hline
\end{tabular}
\end{table}

A Trilha de Pesquisa em SI do SBSI 2026 é o corpus-alvo. Os dois conjuntos do ENIAC 2025 servem como contraste e cumprem funções distintas. A \textit{Main track} pertence a outra comunidade e a outra área, e mantém em comum com o SBSI apenas o formato editorial e o filtro institucional da SBC.

A \textit{Undergraduate track} oferece o contraste mais informativo. Ela compartilha com a \textit{Main track} o evento, a chamada, o período, o corpo de revisores e a área de conhecimento. Varia apenas o subgênero, já que a trilha se destina a trabalhos de graduação e reconhece como contribuição legítima aquilo que a trilha principal não reconheceria. Derivar uma especificação de cada uma e comparar as regras resultantes testa diretamente a premissa que sustenta a escolha da unidade de análise: se as duas especificações forem indistinguíveis, a trilha não é o nível de resolução adequado. Nenhum par de trilhas de eventos diferentes permitiria esse teste, porque nele área, comunidade e trilha variariam ao mesmo tempo, e qualquer diferença observada seria ambígua.

Três trilhas é o número mínimo que permite os dois contrastes de que o desenho depende, entre comunidades distintas e entre trilhas de um mesmo evento, e não um limite do artefato. A especificação de gênero é derivada por trilha, de modo que incluir uma nova trilha consiste em uma nova coleta, sem alteração no escopo doa agente. A versão final do artefato pode então cobrir outras trilhas, outras edições e outras conferências que publiquem seus anais de forma aberta, ampliando o corpus à medida que novos conjuntos forem incorporados.

\subsection{Arquitetura em dois fluxos}
\label{sec:arquitetura-conformidade}

O agente organiza-se em dois fluxos. O \textbf{Fluxo A} executa uma vez por trilha, em lote, e tem por produto a Especificação de Gênero. O \textbf{Fluxo B} executa a cada consulta do autor e a consome. A Figura~\ref{fig:arquitetura} apresenta a organização geral.

No Fluxo A não há verificação de conformidade. Sua responsabilidade é caracterizar o gênero da trilha a partir do corpus, identificando no texto dos artigos aceitos características suficientes para enunciar regras. A conformidade só passa a existir no Fluxo B, quando um manuscrito é confrontado com a especificação já constituída.

\begin{figure}[H]
\centering
\begin{tikzpicture}[
  every node/.style = {font=\footnotesize},
  proc/.style  = {rectangle, draw, rounded corners=2pt, minimum height=9mm,
                  text width=24mm, align=center, inner sep=2pt},
  det/.style   = {proc, fill=black!5},
  inf/.style   = {proc, fill=black!20},
  data/.style  = {rectangle, draw, dashed, minimum height=9mm,
                  text width=24mm, align=center, inner sep=2pt},
  store/.style = {rectangle, draw, dashed, minimum height=9mm,
                  text width=56mm, align=center, inner sep=3pt},
  seta/.style  = {-{Latex[length=2mm]}, thick},
  banda/.style = {font=\footnotesize\itshape, anchor=west}
]

% Fluxo A: construcao da Especificacao de Genero (lote, por trilha)
\node[banda] at (-0.6,1.15) {Fluxo A: por trilha, em lote};

\node[data] (corpus) at (0,0)    {Corpus da trilha\\[-1pt] \scriptsize (artigos aceitos)};
\node[det]  (a1)     at (3.1,0)  {\textbf{A1}\\[-1pt] Ingestão e conversão\\ dos documentos};
\node[inf]  (a2)     at (6.2,0)  {\textbf{A2}\\[-1pt] Extração de\\ \textbf{moves}};
\node[det]  (a3)     at (9.3,0)  {\textbf{A3}\\[-1pt] Agregação\\ determinística};
\node[det]  (a4)     at (12.4,0) {\textbf{A4}\\[-1pt] Derivação das\\ regras de gênero};

\draw[seta] (corpus) -- (a1);
\draw[seta] (a1) -- (a2);
\draw[seta] (a2) -- (a3);
\draw[seta] (a3) -- (a4);

% Artefato de conhecimento intermediario
\node[store] (espec) at (7.75,-2.0)
  {\textbf{Especificação de Gênero}\\[-1pt]
   \scriptsize regras de gênero \,$+$\, segmentos rotulados};

\draw[seta] (a4.south) .. controls +(0,-1.1) and +(1.6,0) .. (espec.east);

% Fluxo B: verificacao (por consulta)
\node[banda] at (-0.6,-5.15) {Fluxo B: por consulta};

\node[data] (manu) at (0,-4)    {Manuscrito\\[-1pt] \scriptsize do autor};
\node[inf]  (b12)  at (3.1,-4)  {\textbf{B1--B2}\\[-1pt] Ingestão e extração\\ \scriptsize (idem A1--A2)};
\node[det]  (b3)   at (6.2,-4)  {\textbf{B3}\\[-1pt] Verificação de\\ conformidade};
\node[inf]  (b4)   at (9.3,-4)  {\textbf{B4}\\[-1pt] Composição\\ do retorno};
\node[data] (feed) at (12.4,-4) {Retorno\\[-1pt] ao autor};

\draw[seta] (manu) -- (b12);
\draw[seta] (b12) -- (b3);
\draw[seta] (b3) -- (b4);
\draw[seta] (b4) -- (feed);

\draw[seta] (espec.south -| b3.north) -- (b3.north);
\draw[seta] (espec.south -| b4.north) -- (b4.north);

% consulta a especificacao, sem manuscrito
\draw[seta, dotted] (espec.west) .. controls +(-2.2,0) and +(0,1.4) .. (b12.north);

% Legenda
\node[det,  minimum height=4mm, text width=3mm] (lg1) at (3.3,-6.1) {};
\node[anchor=west, font=\scriptsize] at (3.6,-6.1) {determinístico};
\node[inf,  minimum height=4mm, text width=3mm] (lg2) at (7.0,-6.1) {};
\node[anchor=west, font=\scriptsize] at (7.3,-6.1) {por inferência (LLM)};
\node[data, minimum height=4mm, text width=3mm] (lg3) at (10.9,-6.1) {};
\node[anchor=west, font=\scriptsize] at (11.2,-6.1) {dado};

\end{tikzpicture}
\caption{Arquitetura do agente de conformidade de gênero. O Fluxo A é executado uma vez por trilha e produz a Especificação de Gênero; o Fluxo B é executado a cada consulta e a consome. A seta pontilhada representa a consulta à especificação sem manuscrito. Blocos claros são determinísticos; blocos escuros envolvem inferência de modelo de linguagem, aplicada de modo idêntico ao corpus (A2) e ao manuscrito (B1--B2). Elaboração própria.}
\label{fig:arquitetura}
\end{figure}

A Tabela~\ref{tab:responsabilidades} detalha o que é determinístico e o que é feito por inferência em cada etapa. Identificar funções retóricas em texto é tarefa de interpretação linguística e cabe ao modelo de linguagem; consolidar o corpus, comparar valores com limiares e decidir sobre conformidade são operações de contagem e cabem a código determinístico.

\begin{table}[H]
\centering
\caption{Distribuição de responsabilidades entre componentes determinísticos e inferenciais}
\label{tab:responsabilidades}
\small
\begin{tabular}{|p{1.6cm}|p{7.6cm}|p{4.0cm}|}
\hline
\textbf{Etapa} & \textbf{Operação} & \textbf{Natureza} \\
\hline
A1 & Ingestão e conversão dos documentos & Determinística \\
\hline
A2 & Extração de \textbf{moves} & Por inferência \\
\hline
A3 & Agregação e atribuição de estatuto & Determinística \\
\hline
A4 & Derivação das regras de gênero & Determinística \\
\hline
B1--B2 & Ingestão e extração do manuscrito & Idem A1--A2 \\
\hline
B3 & Decisão de conformidade & Determinística \\
\hline
B4 & Composição do retorno & Por inferência \\
\hline
\end{tabular}
\end{table}

\subsection{Fluxo A: construção da Especificação de Gênero}
\label{sec:fluxo-a}

O Fluxo A percorre quatro etapas. Na \textbf{ingestão} (A1), os artigos aceitos da trilha, obtidos em PDF, são convertidos em uma representação estruturada da qual se obtêm segmentação de seções, tipagem de itens e contagem de tabelas e figuras. A etapa não envolve inferência de modelo de linguagem, sendo determinística, idempotente e reexecutável sobre o mesmo insumo com resultado idêntico.

Na \textbf{extração de moves} (A2), cada artigo é submetido individualmente a um modelo de linguagem, que rotula os segmentos da introdução segundo o vocabulário do modelo CARS revisado \cite{ResearchGenres} e devolve a saída em um esquema estruturado. Cada chamada opera sobre um contexto reduzido e verificável, o que permite auditar e reprocessar um artigo isolado sem refazer o conjunto. Esta é a única etapa do fluxo em que um modelo de linguagem interpreta o texto.

Na \textbf{agregação} (A3), os registros produzidos em A2 são consolidados em distribuições de frequência da trilha, como a proporção de artigos em que cada \textbf{move} ocorre, a ordem modal de ocorrência e estatísticas estruturais do número de referências, figuras, tabelas e páginas. Nenhuma inferência intervém nesta etapa.

Uma distribuição de frequências apenas descreve. Para que o Fluxo B possa decidir o que merece observação e o que é escolha legítima do autor, cada elemento recebe um \textbf{estatuto}, conforme a classificação de \citeonline{ResearchGenres}, que distingue elementos obrigatórios, opcionais e prováveis em determinadas áreas (PDA). A contribuição desta etapa é resolver empiricamente essa classificação para a trilha analisada: onde Swales indica que o estatuto varia conforme o campo disciplinar, o agente o mede. A atribuição usa a proporção de artigos da trilha em que o elemento ocorre ao menos uma vez, com dois cortes apresentados na Tabela~\ref{tab:estatutos}. Os valores de corte são parâmetros de projeto, calibráveis por trilha.

\begin{table}[H]
\centering
\caption{Estatutos atribuídos aos elementos de gênero}
\label{tab:estatutos}
\small
\begin{tabular}{|p{3.2cm}|p{2.6cm}|p{7.4cm}|}
\hline
\textbf{Estatuto} & \textbf{Ocorrência} & \textbf{Tratamento da ausência} \\
\hline
Obrigatório & $\geq$ 60\% & Ausência é sinalizada \\
\hline
PDA & 30\% a 60\% & Ausência só é observada se a trilha-alvo espera o elemento \\
\hline
Opcional & $<$ 30\% & Ausência não gera observação \\
\hline
\end{tabular}
\end{table}

Na \textbf{derivação das regras} (A4), a distribuição de frequências é traduzida na forma de um predicado. \citeonline{GenresOrgComm} sustentam que gêneros são encenados por meio de regras de gênero, que associam elementos apropriados de forma e substância a situações recorrentes, e que essas regras operam com frequência de modo tácito. O que o Fluxo A realiza é a explicitação dessas regras a partir do único registro em que elas se encontram depositadas, o conjunto de textos que a comunidade reconheceu. Cada regra derivada é um objeto composto, descrito na Tabela~\ref{tab:regra}. Predicado, estatuto e suporte resultam de contagem sobre os registros de A2; a função é o único elemento enunciado por modelo de linguagem, uma vez por regra e revisável antes de qualquer uso.

\begin{table}[H]
\centering
\caption{Composição de uma regra de gênero}
\label{tab:regra}
\small
\begin{tabular}{|p{2.6cm}|p{7.2cm}|p{3.4cm}|}
\hline
\textbf{Elemento} & \textbf{Conteúdo} & \textbf{Origem} \\
\hline
Dimensão & Dimensão do sistema de gêneros a que a regra se refere & \cite{GenreSystems} \\
\hline
Predicado & Condição verificável sobre a representação do manuscrito & A3 \\
\hline
Estatuto & Obrigatória, PDA ou opcional na trilha & Limiares aplicados em A3 \\
\hline
Suporte & Proporção de artigos aceitos que realizam a regra & A3 \\
\hline
Função & Propósito comunicativo que o elemento cumpre no gênero & A4, por inferência \\
\hline
Exemplares & Segmentos do corpus que realizam a regra & A2 \\
\hline
\end{tabular}
\end{table}

A Especificação de Gênero, produto do Fluxo A, organiza-se em três camadas. A \textbf{camada estrutural} reúne as regularidades verificáveis por inspeção, correspondentes às dimensões explicitamente codificadas do sistema de gêneros. A \textbf{camada retórica} contém o modelo de \textbf{moves} da trilha, com o estatuto empírico de cada elemento. A \textbf{camada de exemplares} reúne os segmentos textuais rotulados em A2 que sustentam a demonstração ao autor.

\subsection{Fluxo B: verificação de conformidade}
\label{sec:fluxo-b}

No Fluxo B o autor submete seu manuscrito e indica a trilha-alvo. O agente não pressupõe trilha nem adota valor padrão: sem trilha declarada não há especificação contra a qual verificar, e a operação não se inicia.

O manuscrito percorre primeiro as etapas de \textbf{ingestão e extração} (B1--B2), com o mesmo conversor, o mesmo extrator e o mesmo esquema de saída de A1 e A2. O produto é um registro estruturado na mesma forma que os registros do corpus, condição para que a comparação seguinte seja legítima, pois qualquer viés do extrator incide igualmente sobre os dois lados.

Na \textbf{verificação de conformidade} (B3), cada regra da Especificação de Gênero tem seu predicado avaliado contra o registro do manuscrito. A operação é determinística e produz dois conjuntos, as regras invocadas e as não invocadas. Sobre o segundo aplica-se o filtro de estatuto, de modo que apenas regras de estatuto esperado dão origem a observação, ao passo que a não realização de regras opcionais permanece sem consequência. A etapa não produz pontuação. A recusa de emitir um índice agregado de conformidade é proposital: uma nota converteria retorno formativo em julgamento e aproximaria o agente do papel de revisor. Uma instância pode legitimamente deixar de invocar parte do repertório sem perder pertencimento \cite{GenresOrgComm}, e o agente verifica suficiência, não completude.

Na \textbf{composição do retorno} (B4), cada elemento decidido em B3 é convertido em texto por um modelo de linguagem que recebe no contexto a regra, seu estatuto, seu suporte no corpus, o trecho correspondente do manuscrito e os exemplares associados. Ao modelo cabe a redação, e apenas ela. Não lhe cabe decidir o que está em conformidade, decisão já tomada em B3, nem produzir redação substituta.

O retorno compõe-se de quatro elementos. O \textbf{reconhecimento das regras invocadas}, ou seja, o que o manuscrito realiza e qual função cada realização cumpre. O \textbf{registro das regras esperadas não realizadas}, cada uma acompanhada da função que cumpre no gênero, da proporção de artigos aceitos que a realizam e de exemplares reais que a demonstram. O \textbf{posicionamento descritivo} do manuscrito frente às faixas observadas no corpus, apresentado como contexto comparativo, e não como desvio. E o \textbf{registro informativo} da variação não observada no corpus, isto é, elementos presentes no manuscrito para os quais a especificação não tem correspondente, registrados e nunca corrigidos. A Tabela~\ref{tab:tipologia} sintetiza o tratamento de cada tipo de divergência.

\begin{table}[H]
\centering
\caption{Tipologia das divergências e seu tratamento}
\label{tab:tipologia}
\small
\begin{tabular}{|p{4.0cm}|p{4.6cm}|p{4.8cm}|}
\hline
\textbf{Tipo} & \textbf{Fundamento} & \textbf{Tratamento} \\
\hline
Regra esperada não realizada & Regra distintiva não invocada \cite{GenresOrgComm} & Observação prioritária, com função e exemplares \\
\hline
Realização parcial & \textbf{Move} presente sem os \textit{steps} típicos \cite{ResearchGenres} & Observação com indicação do que a função requer \\
\hline
Desvio de ordem & Sequência divergente da ordem modal da trilha & Observação com ressalva de que a ordem admite variação \\
\hline
Desvio estrutural & Divergência na camada estrutural & Observação objetiva, verificável pelo autor \\
\hline
Regra opcional não realizada & Escolha autoral legítima \cite{ResearchGenres} & Não gera observação \\
\hline
Elemento sem correspondência na especificação & Variação não observada no corpus & Registro informativo, nunca corretivo \\
\hline
\end{tabular}
\end{table}

Tratar a variação não observada como informativa preserva a possibilidade de inovação retórica, que um sistema estritamente normativo penalizaria.


\section{Agente de fronteira de rejeição}
\label{sec:agente2}

A conformidade de gênero é necessária, mas não suficiente. Um artigo pode realizar todas as regras esperadas por uma trilha e ainda ser recusado por razões que a forma não revela, o caso do manuscrito conforme, porém pode ter uma fragilidade no conteúdo. O agente de fronteira de rejeição cobre essa possível fragilidade, ao sinalizar riscos que a conformidade de gênero por construção não alcança.

Sua função é reportar risco de recusa, e não julgar o manuscrito. Ele não afirma que o texto tem falhas, o que reintroduziria o papel de revisor e comprometeria a avaliação do agente de conformidade. O que ele faz é localizar artigos rejeitados semelhantes ao manuscrito, apurar os motivos pelos quais foram recusados e sinalizar ao autor que o manuscrito pode estar exposto aos mesmos motivos. A decisão sobre o que fazer com esse alerta permanece com o autor.

\subsection{Corpus}
\label{sec:corpus-fronteira}

O agente opera sobre um corpus de avaliações de artigos rejeitados. Para o contexto brasileiro não existe corpus público desse tipo, e sua construção fica como trabalho futuro. Nesta primeira versão, o agente usa o ReviewArena \cite{ReviewArena}, base de conferências internacionais que reúne pareceres de artigos aceitos e rejeitados. Os motivos de recusa apurados refletem as comunidades que compõem a base, e não uma trilha brasileira.
 
O ReviewArena é uma escolha da versão atual, condicionada ao que existe publicamente hoje, e não um compromisso do artefato. Suas avaliações provêm do OpenReview e cobrem exclusivamente conferências internacionais de aprendizado de máquina e processamento de linguagem natural, o que as afasta do gênero das conferências brasileiras. Por isso a base pode ser substituída assim que houver alternativa que sirva melhor ao objeto desta pesquisa, seja um corpus de avaliações brasileiras, seja outra base de revisões de conferências mais próxima do domínio de Sistemas de Informação. Como o grafo de conhecimento é construído com um fluxo de processamento a partir da base de entrada, a substituição consiste em reexecutá-lo sobre o novo corpus, sem alteração no restante do agente. 
 
\subsection{Arquitetura em dois fluxos}
\label{sec:arquitetura-fronteira}

O agente segue a mesma organização do agente de conformidade, com um Fluxo A executado uma vez e um Fluxo B executado a cada consulta. A Figura~\ref{fig:arquitetura-fronteira} apresenta os dois momentos.

\subsection{Fluxo A: construção do grafo de conhecimento}
\label{sec:fluxo-a-fronteira}

O Fluxo A constrói, a partir do ReviewArena, um grafo de conhecimento. Cada artigo rejeitado da base é um nó, ligado aos motivos de recusa extraídos de suas avaliações e às características que permitem aproximá-lo de outros artigos. A extração dos motivos de recusa a partir do texto das avaliações é feita por modelo de linguagem, ao passo que a montagem e a indexação do grafo são determinísticas. Esse uso do modelo para derivar o grafo a partir dos documentos de origem, com a consulta posterior feita sobre a estrutura resultante, segue o desenho da abordagem GraphRAG \cite{EdgeGraphRAG2024}. Esse fluxo é executado uma vez e reexecutado apenas quando a base de artigos rejeitados é atualizada.
 
Grafos de conhecimento armazenam conhecimento de forma estruturada e explícita, o que permite fornecê-lo ao modelo de linguagem como base externa de inferência e ao mesmo tempo torna interpretável o caminho pelo qual uma resposta foi obtida \cite{PanKGLLM2024}. E a recuperação orientada por grafo aproveita as relações entre entidades, o que a recuperação vetorial não captura, resultando em recuperação mais precisa e sensível ao contexto \cite{PengGraphRAG}. Isso permite tornar explícitas as relações entre artigo, motivo de recusa e características de aproximação, de modo que o vizinho recuperado seja rastreável e o motivo reportado ao autor possa ser justificado. A definição do que torna dois artigos semelhantes é uma decisão de projeto que condiciona o que o agente reporta: uma aproximação temática recupera rejeitados do mesmo assunto do manuscrito; outras dimensões recuperam outros vizinhos e, com eles, outros motivos de recusa. Essa escolha é registrada como parâmetro de projeto, a ser calibrado quando o grafo for construído.

\subsection{Fluxo B: consulta e retorno de risco}
\label{sec:fluxo-b-fronteira}

O Fluxo B recebe o manuscrito e, como insumo adicional, a análise produzida pelo agente de conformidade. A partir dessas entradas, um modelo de linguagem consulta o grafo em busca de artigos rejeitados semelhantes ao manuscrito. Recuperados os vizinhos, seus motivos de recusa, já gravados no grafo desde o Fluxo A, são reunidos, e o modelo compõe um retorno que apresenta ao autor os riscos a que o manuscrito está exposto e os exemplares de recusa que os fundamentam. A busca e a redação são assistidas pelo modelo, ao passo que os motivos reportados provêm do grafo. A análise do agente de conformidade entra como sinal de aproximação, pois as regras de gênero que o manuscrito realiza ou deixa de realizar orientam a busca por artigos rejeitados comparáveis.

\begin{figure}[H]
\centering
\begin{tikzpicture}[
  every node/.style = {font=\footnotesize},
  proc/.style  = {rectangle, draw, rounded corners=2pt, minimum height=9mm,
                  text width=24mm, align=center, inner sep=2pt},
  det/.style   = {proc, fill=black!5},
  inf/.style   = {proc, fill=black!20},
  data/.style  = {rectangle, draw, dashed, minimum height=9mm,
                  text width=24mm, align=center, inner sep=2pt},
  store/.style = {rectangle, draw, dashed, minimum height=9mm,
                  text width=62mm, align=center, inner sep=3pt},
  seta/.style  = {-{Latex[length=2mm]}, thick},
  banda/.style = {font=\footnotesize\itshape, anchor=west}
]

% Fluxo A: construcao do grafo (uma vez)
\node[banda] at (-0.6,1.15) {Fluxo A: uma vez, sobre o ReviewArena};

\node[data] (base)  at (0,0)    {ReviewArena\\[-1pt] \scriptsize (aceitos e\\ rejeitados)};
\node[inf]  (fa1)   at (3.3,0)  {Extração dos\\ motivos de recusa};
\node[det]  (fa2)   at (6.6,0)  {Montagem e\\ indexação do grafo};

\draw[seta] (base) -- (fa1);
\draw[seta] (fa1) -- (fa2);

% Grafo de conhecimento (artefato intermediario), centrado sob o Fluxo A
\node[store] (kg) at (6.6,-2.0)
  {\textbf{Grafo de conhecimento}\\[-1pt]
   \scriptsize artigos rejeitados, motivos de recusa e relações};

\draw[seta] (fa2.south) -- (kg.north);

% Fluxo B: consulta (por manuscrito)
\node[banda] at (-0.6,-5.15) {Fluxo B: por consulta};

\node[data] (manu)  at (0,-4)    {Manuscrito\\[-1pt] \scriptsize $+$ análise de conformidade};
\node[inf]  (fb1)   at (3.3,-4)  {Busca de artigos\\ rejeitados similares};
\node[det]  (fb2)   at (6.6,-4)  {Leitura dos\\ motivos de recusa};
\node[inf]  (fb3)   at (9.9,-4)  {Composição do\\ retorno de risco};
\node[data] (out)   at (13.0,-4) {Riscos\\[-1pt] ao autor};

\draw[seta] (manu) -- (fb1);
\draw[seta] (fb1) -- (fb2);
\draw[seta] (fb2) -- (fb3);
\draw[seta] (fb3) -- (out);

% o grafo alimenta a busca e a leitura dos motivos: setas verticais retas
\draw[seta] (kg.south -| fb1.north) -- (fb1.north);
\draw[seta] (kg.south -| fb2.north) -- (fb2.north);

% Legenda
\node[det,  minimum height=4mm, text width=3mm] (lg1) at (3.3,-6.1) {};
\node[anchor=west, font=\scriptsize] at (3.6,-6.1) {determinístico};
\node[inf,  minimum height=4mm, text width=3mm] (lg2) at (7.0,-6.1) {};
\node[anchor=west, font=\scriptsize] at (7.3,-6.1) {por inferência (LLM)};
\node[data, minimum height=4mm, text width=3mm] (lg3) at (10.9,-6.1) {};
\node[anchor=west, font=\scriptsize] at (11.2,-6.1) {dado};

\end{tikzpicture}
\caption{Arquitetura do agente de fronteira de rejeição. O Fluxo A constrói o grafo de conhecimento a partir do ReviewArena e é executado uma vez, atualizado apenas quando a base de rejeitados muda; o Fluxo B consulta o grafo a cada manuscrito e recebe, como insumo, a análise do agente de conformidade.}
\label{fig:arquitetura-fronteira}
\end{figure}

\section{Delimitações e riscos de projeto}
\label{sec:limitacoes}

\textbf{Viés de sobrevivência do corpus de aceitos.} O corpus do agente de conformidade é composto apenas por artigos aceitos. Sua especificação modela o gênero tal como realizado pelos textos aprovados, e não a fronteira que separa aceitação de rejeição. Esse agente pode afirmar que um manuscrito diverge do padrão dos aceitos, mas não que tal divergência causaria sua recusa. 
% Cobrir essa fronteira é a função do agente de fronteira de rejeição, cuja limitação atual é distinta: nesta versão ele se apoia no ReviewArena, de comunidades internacionais, de modo que os motivos de recusa que reporta não correspondem ainda a uma trilha brasileira.

\textbf{Dependência da qualidade da extração.} Toda a cadeia do agente de conformidade repousa sobre a acurácia do rotulador de A2. Se a extração erra sistematicamente em um \textbf{move}, as frequências agregadas em A3 e as regras derivadas em A4 estarão erradas, ainda que computadas corretamente. Isso torna necessária uma validação preliminar, que meça a concordância entre os rótulos do modelo e uma anotação manual de referência em amostra reduzida, antes de qualquer processamento em escala.

% \textbf{Anotação por um único anotador.} A anotação de referência é produzida pelo próprio pesquisador, de modo que não há medida de concordância entre anotadores humanos que sirva de teto interpretativo. A alegação sustentável não é que o extrator seja acurado em termos absolutos, mas que concorda com uma anotação de referência em determinada proporção, com erros concentrados em categorias identificáveis.

\textbf{Cobertura da camada retórica.} O modelo CARS descreve a organização retórica de introduções, e não do artigo inteiro. A camada retórica restringe-se, portanto, à introdução, ao passo que a camada estrutural cobre o documento completo. Estender a análise retórica a outras seções exigiria modelos de \textbf{moves} próprios para cada uma, o que fica previsto como trabalho futuro.

\textbf{Uma única edição por trilha.} As especificações de gênero derivam de uma edição de cada evento. Regularidades próprias daquela edição, como um tema em destaque ou uma composição atípica do comitê, não se distinguem de regularidades estáveis da trilha. A separação exigiria séries de múltiplas edições, e a limitação deve ser levada em conta na leitura dos resultados.


\section{Síntese}
\label{sec:sintese-solucao}

O agente de conformidade de gênero converte o conhecimento tácito de uma comunidade em uma especificação verificável, por meio de uma cadeia de operações auditáveis. Ele extrai funções retóricas dos artigos aceitos de uma trilha com um modelo de linguagem, consolida as extrações por agregação determinística, deriva dessa agregação um conjunto de regras explicitadas e, de posse delas, verifica o manuscrito do autor, comunicando tanto o que ele realiza quanto o que deixa de realizar, sempre ancorado em evidência textual e em exemplares reais.

O agente de fronteira de rejeição reporta os motivos recorrentes de recusa a partir de avaliações reais e, nesta versão, opera sobre o ReviewArena. Juntos, os dois agentes cobrem tanto o que caracteriza os artigos aceitos de uma trilha quanto os motivos pelos quais artigos são recusados, recortes que a conformidade de gênero, isoladamente, não alcança. Um corpus de avaliações brasileiras, ainda a ser constituído, tornaria os motivos de recusa específicos da comunidade-alvo.